Note first that it follows from 1.2.10 and 3.3.6 that (i) is a consequence of (ii) and (iii), and that (ii) and (iii) determine the $T_\alpha$ completely. We shall make the construction by induction on $\operatorname{ord}_\Delta(\alpha)$. If $\operatorname{ord}_\Delta(\alpha)=1$, then $\alpha\in\Delta$, and put $T_\alpha=T_i$ if $\alpha=\alpha_i$. Consider the hypothesis:
\begin{quote}
$(H_p)$ there exist $T_\alpha$, for $\alpha\in\mathcal P(\Delta)$, $\operatorname{ord}_\Delta(\alpha)\leq p$, satisfying (ii) and condition (iii) whenever $\operatorname{ord}_\Delta(\alpha)\leq p$, $\operatorname{ord}_\Delta(\beta)\leq p$.
\end{quote}
This holds for $p=1$: indeed, if $\alpha$ and $s_{\alpha_i}(\alpha)=\beta$ are simple, $\alpha_i$ and $\alpha$ are orthogonal; thus, putting $\alpha_j=\alpha=\beta$, one has $n_{ij}=2$, whence
\[
T_iT_jT_i=T_j.
\]
Suppose $p>1$ and $(H_{p-1})$ holds.

\textbf{A)} \emph{Construction of the $T_\alpha$ for $\operatorname{ord}_\Delta(\alpha)\leq p$.} Clearly, it suffices to do this for $\operatorname{ord}_\Delta(\alpha)=p$. There then exists $\alpha_i\in\Delta$ such that $s_{\alpha_i}(\alpha)\in\mathcal P(\Delta)$ and $\operatorname{ord}_\Delta(s_{\alpha_i}(\alpha))<p$ (3.3.3). One will then put
\begin{equation}\tag{$\star$}\label{XXI.5.eq.star}
T_\alpha=T_iT_{s_{\alpha_i}(\alpha)}T_i.
\end{equation}
Let us check that $T_\alpha$ depends only on $\alpha$. Let therefore $\alpha_j\in\Delta$ be such that $s_{\alpha_j}(\alpha)\in\mathcal P(\Delta)$ and $\operatorname{ord}_\Delta(s_{\alpha_j}(\alpha))<p$. Let us prove that
\begin{equation}\tag{+}\label{XXI.5.eq.plus}
T_iT_{s_{\alpha_i}(\alpha)}T_i=T_jT_{s_{\alpha_j}(\alpha)}T_j.
\end{equation}
Distinguish two cases.

(1) Suppose $\alpha$ a linear combination of $\alpha_i$ and $\alpha_j$. The same is then true of $s_{\alpha_i}(\alpha)$ and $s_{\alpha_j}(\alpha)$, and, by $(H_{p-1})$, $T_{s_{\alpha_i}(\alpha)}$ and $T_{s_{\alpha_j}(\alpha)}$ are written as words in $T_i$ and $T_j$. Since the projection of (+) into $W$ holds, and the theorem is true for $n=2$, hence $p$ is injective on the subgroup of $\overline W$ generated by $T_i$ and $T_j$, (+) indeed holds.

(2) Suppose $\alpha$ is not a linear combination of $\alpha_i$ and $\alpha_j$. Then if $w\in W_{\alpha_i,\alpha_j}$, the $w(\alpha)$ are all positive (cf. 3.4.9). The relation to be checked is also written
\begin{equation}\tag{++}\label{XXI.5.eq.plusplus}
(T_iT_j)^{n_{ij}-1}T_{s_{\alpha_i}(\alpha)}(T_jT_i)^{n_{ij}-1}=T_{s_{\alpha_j}(\alpha)}.
\end{equation}
Now it follows from 4.4 that the $w(\alpha)$ all have order $<p$ for $w\in W_{\alpha_i,\alpha_j}$, $w\ne1$. One may therefore apply hypothesis $(H_{p-1})$ $2(n_{ij}-1)$ times, and the proof is finished.

\textbf{B)} \emph{Verification of $(H_p)$.}{\renewcommand{\thefootnote}{(22)}\nde{The original has been expanded in what follows.}} One must check that if $\alpha_j\in\Delta$ and $\beta=s_{\alpha_j}(\alpha)$ satisfies $\operatorname{ord}_\Delta(\beta)\leq p$, then $T_jT_\alpha T_j=T_\beta$. If $\operatorname{ord}_\Delta(\beta)<p$, this follows from what precedes (since $T_j^2=1$); thus one may suppose $\operatorname{ord}_\Delta(\beta)=p=\operatorname{ord}_\Delta(\alpha)$. In this case, $\alpha$ and $\alpha_j$ are orthogonal, hence $\beta=\alpha$, and one must see that
\begin{equation}\tag{$\dagger$}\label{XXI.5.eq.dagger}
T_jT_\alpha T_j=T_\alpha.
\end{equation}
By ($\star$) above, one has $T_\alpha=T_iT_{s_{\alpha_i}(\alpha)}T_i$, and one therefore need only check the following equality:
\begin{equation}\tag{+++}\label{XXI.5.eq.plusplusplus}
T_jT_iT_{s_{\alpha_i}(\alpha)}T_iT_j=T_\alpha=T_iT_{s_{\alpha_i}(\alpha)}T_i.
\end{equation}
Put $m=n_{ij}$ and $s_i=s_{\alpha_i}$, $s_j=s_{\alpha_j}$. One has $T_jT_i=(T_iT_j)^{m-1}$ and, by 4.4, $\operatorname{ord}_\Delta(w(\alpha))<p$ for every $w\in W_{ij}$ distinct from $1=(s_is_j)^m$ and from $s_j=s_i(s_is_j)^{m-1}$. One deduces, by the induction hypothesis, that
\[
\begin{split}
T_j(T_iT_j)^{m-2}T_{s_i(\alpha)}(T_jT_i)^{m-2}T_j
&=T_{s_j(s_is_j)^{m-2}s_i(\alpha)}\\
&=T_{s_is_j(\alpha)}=T_{s_i(\alpha)}
\end{split}
\]
(the last equality since $s_j(\alpha)=\alpha$), whence finally
\[
(T_iT_j)^{m-1}T_{s_i(\alpha)}(T_jT_i)^{m-1}=T_iT_{s_i(\alpha)}T_i,
\]
which proves (+++).

\begin{lemma}{5.3}\label{XXI.5.3}
Let $h\in\overline W$. Write it
\[
h=T_{\alpha_1}\cdots T_{\alpha_m}
\]
with $\alpha_i\in\Delta$, not necessarily distinct, in such a way that $m$ is minimal. Then
\[
p(h)(\alpha_m)\in-\mathcal P(\Delta).
\]
\end{lemma}
Indeed, since $p(T_{\alpha_m})(\alpha_m)=s_{\alpha_m}(\alpha_m)=-\alpha_m$, if $p(h)(\alpha_m)$ were positive, there would exist an index $k$, $1\leq k\leq m-1$, such that
\[
u=s_{\alpha_{k+1}}\cdots s_{\alpha_m}(\alpha_m)=-s_{\alpha_{k+1}}\cdots s_{\alpha_{m-1}}(\alpha_m)\in-\mathcal P(\Delta),
\]
and $s_{\alpha_k}(u)\in\mathcal P(\Delta)$. But then necessarily $u=-\alpha_k$ (3.3.1), whence
\[
s_{\alpha_{k+1}}\cdots s_{\alpha_{m-1}}(\alpha_m)=\alpha_k,
\]
which implies by (iii)
\[
T_{\alpha_k}T_{\alpha_{k+1}}\cdots T_{\alpha_{m-1}}T_{\alpha_m}=T_{\alpha_{k+1}}\cdots T_{\alpha_{m-1}},
\]
and this contradicts the minimality of $m$.

Let now $h\in\overline W$ be such that $p(h)(\mathcal P(\Delta))\subset\mathcal P(\Delta)$. By lemma~5.3, one has $p(h)=1$, which proves theorem~5.1 and, in addition:
\begin{corollary}{5.4}\label{XXI.5.4}
If $R_+$ is a system of positive roots and $w\in W$ is such that $w(R_+)=R_+$, then $w=1$.
\end{corollary}
\begin{corollary}{5.5}\label{XXI.5.5}
The Weyl group acts freely and transitively on the set of systems of positive roots (resp. systems of simple roots, resp. Weyl chambers).
\end{corollary}
Now choose a system of simple roots $\Delta$. Put $R^+=\mathcal P(\Delta)$.

{\renewcommand{\thefootnote}{(23)}\nde{$R^+$ has been written here instead of $R_+$, so as to be able to write $R_{\alpha,\beta}^+=R_{\alpha,\beta}\cap R^+$.}}%
For every pair of simple roots $(\alpha,\beta)\in\Delta\times\Delta$, denote by $R_{\alpha,\beta}$ the set of roots which are linear combinations of $\alpha$ and $\beta$. Put $R_{\alpha,\beta}^+=R^+\cap R_{\alpha,\beta}$, and let $W_{\alpha,\beta}$ be the Weyl group of $R_{\alpha,\beta}$, that is, the subgroup of $W$ generated by $s_\alpha$ and $s_\beta$.
\begin{theorem}{5.6 (Tits)}\label{XXI.5.6}
Let $\alpha$ and $\beta$ be two simple roots, and $w\in W$ such that $w(\alpha)=\beta$. There exist a sequence of simple roots $\alpha_0,\ldots,\alpha_m$ and a sequence $w_0,\ldots,w_{m-1}$ of elements of $W$ satisfying the following conditions:
\begin{enumeratei}
\item $\alpha_0=\alpha$, $\alpha_m=\beta$.
\item $w=w_{m-1}w_{m-2}\cdots w_0$.
\item $w_i(\alpha_i)=\alpha_{i+1}$ for $0\leq i\leq m-1$.
\item For every $i$, $0\leq i\leq m-1$, such that $\alpha_i\ne\alpha_{i+1}$, one has $w_i\in W_{\alpha_i,\alpha_{i+1}}$.
\item For every $i$, $0\leq1\leq m-1$, such that $\alpha_i=\alpha_{i+1}$, there exists a simple root $\beta_i$ such that $w_i\in W_{\alpha_i,\beta_i}$.
\end{enumeratei}
% typo?: 0 <= 1 <= m-1 in (v) kept as printed.
\end{theorem}
Put{\renewcommand{\thefootnote}{(24)}\nde{In what follows, the original has been expanded, and equality (1) corrected.}}
\[
\mathrm M(w)=\operatorname{Card}(R^+\cap w^{-1}(-R^+))=\operatorname{Card}\{\alpha\in R^+\mid w(\alpha)\in-R^+\}.
\]
If $\mathrm M(w)=0$, then $w(R^+)=R^+$, hence $w=1$ by 5.4, and the theorem is trivial ($m=0$, assertions (iii) to (v) are empty). Argue by induction on $\mathrm M(w)$. If $\mathrm M(w)>0$, there exists $\beta_0\in\Delta$ such that $w(\beta_0)\in-R^+$. Put $\alpha_0=\alpha$. Consider the set
\[
A=w^{-1}(R^+)\cap R_{\alpha_0,\beta_0}.
\]
It is a system of positive roots of $R_{\alpha_0,\beta_0}$. There therefore exists $w_0\in W_{\alpha_0,\beta_0}$ such that
\[
w_0^{-1}(R_{\alpha_0,\beta_0}^+)=A.
\]
Put $w'=ww_0^{-1}$. By 3.4.9, one has at once
\[
R^+-R_{\alpha_0,\beta_0}^+=w_0(R^+-R_{\alpha_0,\beta_0}^+),
\]
whence
\begin{equation}\tag{1}\label{XXI.5.6.eq.1}
(R^+-R_{\alpha_0,\beta_0}^+)\cap w^{-1}(-R^+)
=w_0\bigl((R^+-R_{\alpha_0,\beta_0}^+)\cap w'^{-1}(-R^+)\bigr).
\end{equation}
On the other hand,
\begin{equation}\tag{2}\label{XXI.5.6.eq.2}
\beta_0\in R_{\alpha_0,\beta_0}^+\cap w^{-1}(-R^+),
\end{equation}
and, since $w_0(R_{\alpha_0,\beta_0})=R_{\alpha_0,\beta_0}$, one has $w_0(-A)=R_{\alpha_0,\beta_0}\cap w'^{-1}(-R^+)$, whence
\begin{equation}\tag{$2'$}\label{XXI.5.6.eq.2prime}
R_{\alpha_0,\beta_0}^+\cap w'^{-1}(-R^+)=R_{\alpha_0,\beta_0}^+\cap w_0(-A)
=R_{\alpha_0,\beta_0}^+\cap-R_{\alpha_0,\beta_0}^+=\varnothing.
\end{equation}
It follows from (1), (2) and ($2'$) that $\mathrm M(w')<\mathrm M(w)$.

Put $\alpha_1=w_0(\alpha_0)$; let us show that $\alpha_1\in\Delta$, i.e. $\alpha_0\in w_0^{-1}(\Delta)$. One knows that $w(\alpha_0)\in\Delta$, hence $\alpha_0\in w^{-1}(\Delta)$, hence also $\alpha_0\in w^{-1}(\Delta)\cap R_{\alpha_0,\beta_0}$, hence $\alpha_0$ is a simple root of $A=w^{-1}(R^+)\cap R_{\alpha_0,\beta_0}=w_0^{-1}(R_{\alpha_0,\beta_0}^+)$, hence belongs to
\[
w_0^{-1}(\Delta\cap R_{\alpha_0,\beta_0}^+)=w_0^{-1}(\{\alpha_0,\beta_0\})
\]
(see 3.4.8). Thus $\alpha_1=w_0(\alpha_0)$ equals $\alpha_0$ or $\beta_0$. If $\alpha_1\ne\alpha_0$, one has $\alpha_1=\beta_0$ and $w_0\in W_{\alpha_0,\alpha_1}$.

Finally, one has $\beta=w'(\alpha_1)$, with $\mathrm M(w')<\mathrm M(w)$, and one concludes by induction.

\sgasection{6}{Morphisms of root data}\label{XXI.sec.6}
\numpara{6.1} \textbf{Definition.} Let $\mathcal R=(M,M^*,R,R)$ and $\mathcal R'=(M',M'^*,R',R'^*)$ be two root data. Let $f:M'\to M$ be a linear map and ${}^tf:M^*\to M'^*$ its transpose.
% typo?: repeated R rather than R star in the first tuple kept as printed.
\begin{definition}{6.1.1}\label{XXI.6.1.1}
One says that $f$ is a \emph{morphism} from $\mathcal R'$ to $\mathcal R$, and writes $f:\mathcal R'\to\mathcal R$, if $f$ induces a bijection from $R'$ onto $R$ and ${}^tf$ a bijection from $R^*$ onto $R'^*$.
\end{definition}
Then ${}^tf$ is a morphism of the dual root data:
\[
{}^tf:\mathcal R^*\to\mathcal R'^*.
\]
One easily sees that if $f$ is a morphism from $\mathcal R'$ to $\mathcal R$, and one puts $\alpha=f(\alpha')$ for $\alpha'\in R'$, then $\alpha'^*={}^tf(\alpha^*)$. Indeed, one sees immediately that if $p$ and $p'$ denote the maps of 1.2.1 relative to $\mathcal R$ and $\mathcal R'$, respectively, one has $p'={}^tf\circ p\circ f$, and the desired assertion follows at once. We leave to the reader the proofs of the following statements, almost all of which are trivial.
\begin{proposition}{6.1.2}\label{XXI.6.1.2}
Let $f:\mathcal R'\to\mathcal R$ be a morphism of root data. If $\alpha'\in R'$ and $\alpha=f(\alpha')$, then $\alpha'^*={}^tf(\alpha^*)$. Moreover, $f$ induces isomorphisms:
\[
R'\isomto R,\qquad\Gamma_0(R')\isomto\Gamma_0(R),\qquad\mathcal V(R')\isomto\mathcal V(R),
\]
and ${}^tf$ induces isomorphisms:
\[
R^*\isomto R'^*,\qquad\Gamma_0(R^*)\isomto\Gamma_0(R'^*),\qquad\mathcal V(R^*)\isomto\mathcal V(R'^*),
\]
the last being the transpose of the corresponding morphism induced by $f$. The map $s_{\alpha'}\mto s_{f(\alpha')}$ extends to an isomorphism $W(\mathcal R')\isomto W(\mathcal R)$ compatible with the actions of these two groups on the sets of 1.1.13.
\end{proposition}
\begin{proposition}{6.1.3}\label{XXI.6.1.3}
The maps
\[
\Delta'\mto f(\Delta'),\qquad R'_+\mto f(R'_+),\qquad C'\mto({}^tf\otimes\RR)^{-1}(C')
\]
define bijective correspondences between systems of simple roots, systems of positive roots, and Weyl chambers for $\mathcal R'$ and $\mathcal R$. These correspondences are compatible with the action of the Weyl groups and with the correspondences
\[
\mathcal S(R_+)\leftrightarrow R_+\leftrightarrow\mathcal C(R_+).
\]
\end{proposition}
\begin{lemma}{6.1.4}\label{XXI.6.1.4}
Morphisms compose. For the morphism $f:\mathcal R'\to\mathcal R$ to be an isomorphism, it is necessary and sufficient that $f:M'\to M$ be bijective.
\end{lemma}

\sgasection{6.2}{Isogenies}
\begin{definition}{6.2.1}\label{XXI.6.2.1}
A morphism $f:\mathcal R'\to\mathcal R$ of root data is called an \emph{isogeny} if $f:M'\to M$ is injective with finite cokernel.
\end{definition}
If $f$ is an isogeny, ${}^tf$ is an isogeny.
\begin{definition}{6.2.2}\label{XXI.6.2.2}
Let $f:\mathcal R'\to\mathcal R$ be an isogeny. Put $K(f)=\Coker(M'\xrightarrow{f}M)$.
\end{definition}
\begin{lemma}{6.2.3}\label{XXI.6.2.3}
One has a natural pairing
\[
K(f)\times K({}^tf)\to\QQ/\ZZ,
\]
putting these two finite groups in duality.
\end{lemma}
This is classical.
\begin{lemma}{6.2.4}\label{XXI.6.2.4}
If $f:\mathcal R'\to\mathcal R$ is a morphism, then $\operatorname{rgss}(\mathcal R')=\operatorname{rgss}(\mathcal R)$. If in addition $f$ is an isogeny, one also has $\operatorname{rgred}(\mathcal R')=\operatorname{rgred}(\mathcal R)$.
\end{lemma}
Trivial.
\begin{lemma}{6.2.5}\label{XXI.6.2.5}
Every morphism of semisimple root data is an isogeny.
\end{lemma}
This follows at once from the fact that $f$ must induce an isomorphism from $V'=\mathcal V(R')$ onto $V=\mathcal V(R)$.

If $\mathcal R'$ and $\mathcal R$ are semisimple, every isogeny $f:\mathcal R'\to\mathcal R$ defines a commutative diagram:
\[
\xymatrix{\Gamma_0(R')\ar[r]^{\sim}\ar@{^{(}->}[d]&\Gamma_0(R)\ar@{^{(}->}[d]\\
M'\ar@{^{(}->}[r]^f&M.}
\]
If $M=\Gamma_0(R)$, $f$ is necessarily an isomorphism.
\begin{definition}{6.2.6}\label{XXI.6.2.6}
A root datum is called \emph{adjoint} (resp. \emph{simply connected}) if $M=\Gamma_0(R)$, resp. $M^*=\Gamma_0(R^*)$.
\end{definition}
An adjoint or simply connected root datum is therefore semisimple. On the other hand, $\mathcal R$ is adjoint (resp. simply connected) if and only if $\mathcal R^*$ is simply connected (resp. adjoint). By the preceding result, one has:
\begin{proposition}{6.2.7}\label{XXI.6.2.7}
Let $\mathcal R$ be a semisimple root datum. The following conditions are equivalent:
\begin{enumeratei}
\item $\mathcal R$ is adjoint (resp. simply connected).
\item Every isogeny $\mathcal R'\to\mathcal R$ (resp. $\mathcal R\to\mathcal R'$) is an isomorphism.
\end{enumeratei}
\end{proposition}
\begin{proposition}{6.2.8}\label{XXI.6.2.8}
Let $\mathcal R$ be an adjoint (resp. simply connected) root datum. Every indivisible root (resp. coroot) is an indivisible element of $M$ (resp. $M^*$).
\end{proposition}
Indeed, every indivisible root belongs to a basis of $\Gamma_0(R)$ by 3.3.5.

\sgasection{6.3}{Radical and coradical}
Let $\mathcal R$ be a root datum. Put
\[
N=\{x\in M\mid(\alpha^*,x)=0\text{ for every }\alpha^*\in R^*\};\qquad
N^*=M^*/\mathcal V(R^*)\cap M^*.
\]
\begin{lemma}{6.3.1}\label{XXI.6.3.1}
Consider the canonical morphisms:
\[
N\to M,\qquad M^*\to N^*.
\]
They are transposes of one another, and $N^*$ is identified with the dual of $N$.
\end{lemma}
This is immediate, taking account of 1.2.5.
\begin{definition}{6.3.2}\label{XXI.6.3.2}
One calls the \emph{coradical} of $\mathcal R$, and denotes by $\operatorname{corad}(\mathcal R)$, the trivial root datum
\[
\operatorname{corad}(\mathcal R)=(N,N^*,\varnothing,\varnothing).
\]
\end{definition}
If one puts $\mathcal R^0=(M,M^*,\varnothing,\varnothing)$ (this is a trivial root datum), one therefore has a morphism
\[
\operatorname{corad}(\mathcal R)\to\mathcal R^0.
\]
\begin{definition}{6.3.3}\label{XXI.6.3.3}
One calls the \emph{radical} of $\mathcal R$, and denotes by $\operatorname{rad}(\mathcal R)$, the trivial root datum:
\[
\operatorname{rad}(\mathcal R)=\operatorname{corad}(\mathcal R^*)^*.
\]
{\renewcommand{\thefootnote}{(25)}\nde{I.e. if one puts $P=\{y\in M^*\mid(y,\alpha)=0\text{ for every }\alpha\in R\}$ and $P^*=M/\mathcal V(R)\cap M$, one has $\operatorname{rad}(\mathcal R)=(P^*,P,\varnothing,\varnothing)$.}}
\end{definition}
One therefore has a diagram
\[
\xymatrix{\operatorname{corad}(\mathcal R)\ar[rr]\ar[dr]&&\operatorname{rad}(\mathcal R)\\
&\mathcal R^0\ar[ur]&},
\]
whose transpose is the corresponding diagram for $\mathcal R^*$.
\begin{lemma}{6.3.4}\label{XXI.6.3.4}
The canonical morphism $u:\operatorname{corad}(\mathcal R)\to\operatorname{rad}(\mathcal R)$ is an isogeny.
\end{lemma}
\begin{definition}{6.3.5}\label{XXI.6.3.5}
Put $N(\mathcal R)=K(u)=M/((\mathcal V(R)\cap M)+\bigcap_{\alpha\in R}\Ker(\alpha^*))$. One then has a canonical pairing
\[
N(\mathcal R)\times N(\mathcal R^*)\to\QQ/\ZZ.
\]
\end{definition}
\begin{lemma}{6.3.6}\label{XXI.6.3.6}
One has $\operatorname{rgred}(\operatorname{rad}(\mathcal R))=\operatorname{rgred}(\operatorname{corad}(\mathcal R))=\operatorname{rgred}(\mathcal R)-\operatorname{rgss}(\mathcal R)$, and the following conditions are equivalent:
\begin{enumeratei}
\item $\mathcal R$ is semisimple,
\item $\operatorname{rad}(\mathcal R)=0$,
\item $\operatorname{corad}(\mathcal R)=0$.
\end{enumeratei}
\end{lemma}
